\section{Sisteme de Recomandare: Paradigme și Evoluție}
\label{sec:recsys_paradigms}

Sistemele de recomandare (Recommender Systems - RS) au devenit instrumente omniprezente în mediul digital, fundamentând experiența utilizatorului pe platforme majore precum Netflix, Amazon sau Spotify. În domeniul modei, acestea îndeplinesc o funcție critică: filtrarea unui spațiu vast de produse pentru a identifica articolele relevante pentru un utilizator specific.

Literatura de specialitate clasifică sistemele de recomandare în trei categorii principale \cite{adomavicius2005toward}:

\subsection{Filtrarea Colaborativă (Collaborative Filtering)}
Această abordare se bazează pe premisa că utilizatorii care au avut preferințe similare în trecut vor avea preferințe similare și în viitor.
\begin{itemize}
    \item \textbf{User-Based}: Identifică utilizatori cu comportament similar ("Vecini") și recomandă articole apreciate de aceștia.
    \item \textbf{Item-Based}: Calculează similaritatea între articole pe baza co-apariției în preferințele utilizatorilor \cite{su2009survey}.
\end{itemize}
\textit{Limitări:} Problema "Cold Start" (dificultatea de a face recomandări pentru utilizatori noi sau produse noi fără istoric) este un obstacol major în modă, unde colecțiile se schimbă sezonier.

\subsection{Filtrarea Bazată pe Conținut (Content-Based Filtering)}
Această metodă recomandă articole similare cu cele pe care utilizatorul le-a apreciat anterior, bazându-se pe atributele intrinseci ale produselor. În contextul modei, atributele pot fi textuale (brand, material, preț) sau vizuale (culoare, textură, formă) \cite{mcauley2015image}.
\textit{Avantaje:} Rezolvă parțial problema Cold Start pentru produsele noi și permite o explicabilitate mai mare a recomandărilor ("Ți-am recomandat X pentru că este similar cu Y").

\subsection{Sisteme Hibride}
Pentru a depăși limitările individuale ale metodelor de mai sus, sistemele moderne adoptă arhitecturi hibride. O abordare populară este utilizarea rețelelor neurale profunde (Deep Learning) pentru a extrage caracteristici latente din imagini (Content-based) și a le combina cu vectori de embedding ai utilizatorilor \cite{he2016deep}.

\section{Analiza Soluțiilor Existente pe Piață}
Piața aplicațiilor de "Virtual Styling" include câteva soluții notabile, însă majoritatea se concentrează pe segmentul feminin sau pe e-commerce pur:

\begin{enumerate}
    \item \textbf{Stitch Fix}: Un serviciu hibrid care combină algoritmi de recomandare cu stiliști umani. Deși foarte precis, modelul lor de business este scump și greu de scalat.
    \item \textbf{Google Lens / Pinterest Lens}: Oferă funcționalitatea de "Visual Search" (găsește produse similare vizual), dar nu generează ținute complete (\textit{outfits}), ci doar produse individuale.
    \item \textbf{Acloset / Smart Closet}: Aplicații de digitalizare a garderobei care permit crearea manuală de colaje. Automatizarea este limitată, bazându-se pe reguli simpliste (ex: "Random Shuffle").
\end{enumerate}

StyleGenAI se diferențiază prin focusul pe \textbf{generarea automată de ținute armonice cromatic}, eliminând efortul manual al utilizatorului de a potrivi hainele.
